\documentclass[10pt,letterpaper]{article}
\usepackage{url}
\usepackage[colorlinks,urlcolor=blue]{hyperref}
\usepackage[UTF8]{inputenc}
%\usepackage[latin1]{inputenc}
\usepackage[T1]{fontenc}
\usepackage{uarial}
\renewcommand\familydefault{\sfdefault}
\usepackage{amsmath,amssymb,amsfonts,latexsym,cancel, wrapfig, multicol,enumerate,graphicx, xcolor,paralist,tikz,tkz-tab,dsfont}
\usepackage[letterpaper,showframe=false]{geometry}
\usepackage[spanish]{babel}
\usepackage{color}
\usepackage{hyperref}
%colores------------------------------------------------
\definecolor{encabezado}{RGB}{64,107,204} % % % Principal
\definecolor{titulos}{RGB}{112,90,45}
\definecolor{actividades}{RGB}{112,90,45}
\definecolor{cuadro1}{RGB}{204,219,226}
\definecolor{cuadro2}{RGB}{226,210,178} % % % tercero
\definecolor{cuadro3}{RGB}{255,230,179}
\definecolor{CRecapitulacion}{RGB}{204,219,226}
\definecolor{CcontenidosObjetivos}{RGB}{204,219,226} % % Secundario
\definecolor{ContenidosObjetivos}{RGB}{112,90,45}
\definecolor{numerosTabla}{RGB}{0,0,0}


%--------------- Colores IS22-----------------------
%\definecolor{encabezado}{RGB}{0,134,75} % % % Principal
%\definecolor{titulos}{RGB}{1,110,111}
%\definecolor{actividades}{RGB}{1,110,111}
%\definecolor{cuadro1}{RGB}{173,211,144}
%\definecolor{cuadro2}{RGB}{249,241,133} % % % tercero
%\definecolor{cuadro3}{RGB}{253,199,117}
%\definecolor{CRecapitulacion}{RGB}{188,228,230}
%\definecolor{CcontenidosObjetivos}{RGB}{188,228,230} % % Secundario
%\definecolor{ContenidosObjetivos}{RGB}{1,110,111}
%\definecolor{numerosTabla}{RGB}{70,70,70}

%--------------- Colores antes de IS22-----------------------
%\definecolor{encabezado}{RGB}{65,130,185}
%\definecolor{titulos}{RGB}{65,130,185}
%\definecolor{actividades}{RGB}{239,80,0}
%\definecolor{cuadro1}{RGB}{216,216,216}
%\definecolor{cuadro2}{RGB}{249,206,17}
%\definecolor{cuadro3}{RGB}{239,159,77}
%\definecolor{CRecapitulacion}{RGB}{216,216,216}
%\definecolor{CcontenidosObjetivos}{RGB}{138,186,212}
%\definecolor{ContenidosObjetivos}{RGB}{86,86,86}
%\definecolor{numerosTabla}{RGB}{48,103,133}
%-------------------------------------------------------

\DeclareGraphicsExtensions {.pdf,.png,.jpg}

%Dimensiones-------------------------------------------
\textheight=26cm
\textwidth=20 cm
\topmargin=-3.2 cm
\oddsidemargin=-2 cm
\evensidemargin=0 cm
\parskip=0 mm
\parindent=0 mm
%-----------------------------------------------------


%\usepackage[T1]{fontenc}
\usepackage{uarial}
\renewcommand\familydefault{\sfdefault}
%\usepackage{blindtext}

\newcommand{\URLLibroMN}{\href{https://ezproxy.itcr.ac.cr/login?url=https://www.bibliotecaebook.com/my-library}{\textbf{Acceso al e-book}} Libro de Métodos Numéricos. Autor Jeffry Chavarría Molina.}

\newcommand{\CursoOctave}{\href{https://lugfi.github.io/curso-octave/index.html}{\textbf{Ingreso al curso}}}

\newcommand{\PaqueteSymbolic}{\href{https://docs.sympy.org/latest/tutorials/intro-tutorial/index.html}{\textbf{Paquete Sympy:}} Paquete para realizar cálculo simbólico en Python}
%--------------------------------
%Videos


\newcommand{\URLVideoUno}{\href{https://youtu.be/ajYx1_p9ENk?si=HcU6I23R4iOyACXt}{\textbf{Vídeo 1:}} E. D. con Métodos Numéricos Teoría}

\newcommand{\URLVideoDos}{\href{https://drive.google.com/file/d/16GAE1hLeBJZZ_wsUSGfVm6QBWMIPFX5s/view}{\textbf{Video 2:}} Introducción a E. D. Numéricas}

\newcommand{\URLVideoTres}{\href{https://youtu.be/-z-xH9WdpSs}{\textbf{Video 3:}} ED01: Método de Euler}

\newcommand{\URLVideoCuatro}{\href{https://drive.google.com/file/d/1cLupAmFHfxKZ5KFbxQljqTgPBWySScXh/view}{\textbf{Video 4:}} Euler}

\newcommand{\URLVideoCinco}{\href{https://youtu.be/AmYC7WBAZPg}{\textbf{Video 5:}} Parte 1}
\newcommand{\URLVideoSeis}{\href{https://youtu.be/U32K6xWgmoA}{\textbf{Video 6:}} Parte 2}



\newcommand{\URLVideoSiete}{\href{https://youtu.be/5K-DbO3tiLU}{\textbf{Video 7:}} Método de Taylor}
\newcommand{\URLVideoOcho}{\href{https://youtu.be/sof7m9UT8oU}{\textbf{Video 8:}} ED02: Método de Taylor de orden n para EDO}

\newcommand{\URLVideoDiez}{\href{https://youtu.be/8PIiLIDGE1k}{\textbf{Video 9:}} ED03: Método de Euler mejorado}
\newcommand{\URLVideoNueve}{\href{https://youtu.be/VDL00Fs0AxE}{\textbf{Video 9:}} Programación Método de Taylor en Octave }




%------------Vídeos teoricos---------------
\newcommand{\URLVideoTeoriaActiI}{\href{https://drive.google.com/file/d/1Is-lhFgiC0lT138Ikfc-pgqc2trrvBBq/view?usp=sharing}{\textbf{Vídeo de teoría 1:}}  Complementario al libro}
\newcommand{\URLVideoTeoriaActiII}{\href{https://drive.google.com/file/d/1kloExvEhsScDhAK6WMXRjhp5rQ62zYDm/view?usp=sharing}{\textbf{Vídeo de teoría 2:}} Complementario al libro}
\newcommand{\URLVideoTeoriaActiIII}{\href{https://drive.google.com/file/d/1H7o1l6RkJaNmG_7-_cLjfnSolojOus38/view?usp=sharing}{\textbf{Vídeo de teoría 3:}} Complementario al libro}


%---------Infografías-----------------



%----------Solucionario-------------
\newcommand{\FolletoPractica}{\href{https://tecnube1-my.sharepoint.com/:b:/g/personal/jchavarria_itcr_ac_cr/ESbC_k0vPWdMovfv0-d9GpoBcM31nPWa26ScY419frpBdw?e=hutomi}{\textbf{Folleto de práctica}} Folleto completo de práctica del curso.}

\newcommand{\FolletoPracticaDOS}{\href{https://tecnube1-my.sharepoint.com/:b:/g/personal/jchavarria_itcr_ac_cr/ESbC_k0vPWdMovfv0-d9GpoBcM31nPWa26ScY419frpBdw?e=hutomi}{\textbf{Folleto de práctica}}}


%----------Podcast---------------------------------------------

\newcommand{\URLpodcast}{\href{https://soundcloud.com/angie-solis-palma/logica-matematica-y-conectivas-logicas?utm_source=clipboard\&utm_campaign=wtshare\&utm_medium=widget\&utm_content=https\%253A\%252F\%252Fsoundcloud.com\%252Fangie-252Fsol-palmamama}{\textbf{Podcast:}} Lógica matemática.}



\newcommand{\URLFacebook}{\href{https://tecdigital.tec.ac.cr/dotlrn/catedras.tec/matematicadiscreta/one-community?page_num=0&loc=}{enlace}}

\newcommand{\URLPortalComunidad}{\href{https://tecdigital.tec.ac.cr/dotlrn/catedras.tec/matematicadiscreta/one-community?page_num=0&loc=}{portal de la comunidad del TecDigital para MD}}

\newcommand{\URLPortalGrupoTecDigital}{\href{https://tecdigital.tec.ac.cr/dotlrn/classes/MA/MA1403/S-2-2022.CA.MA1403.6/one-community?page_num=0&loc=}{portal del TecDigital para el grupo 06}}

\newcommand{\URLConsulta}{\href{https://tecdigital.tec.ac.cr/dotlrn/catedras.tec/matematicadiscreta/one-community?page_num=5&loc=}{enlace}}

\newcommand{\ImagenDM}{\includegraphics[scale=0.7]{imagenes/MN.png}}

\newcommand{\Motivacion}{\includegraphics[scale=0.6]{imagenes/motivacion.png}}


\begin{document}
\setlength{\parskip}{0.3cm}

%Encabezado-----------------------------------------------------------------------------------	






\begin{center}
{ \fboxsep 12pt
	\fcolorbox {encabezado}{encabezado}{
	\begin{minipage}[t]{19.5cm}
		{\color{white} 
		TEC | MÉTODOS NUMÉRICOS GRUPO SEMIPRESENCIAL \hfill  II SEMESTRE 2026
 \\
		\rule[0.1cm]{19.5cm}{0.01cm}
		\begin{center}
			{\Huge  \textbf{Semana 13}}
		\end{center}
		\vspace{-0.5cm}
		\rule[0cm]{19.5cm}{0.01cm}
		\begin{center}
			{  \textbf{Solución numérica de Ecuaciones Diferenciales}}
		\end{center}
					}
	\vspace{-0.5cm}
	\end{minipage}
	} 
}
\end{center}


%Pagina1-------------------------------------------------------------------------
\begin{minipage}[c]{14cm}
	\begin{center}
		\includegraphics[height=5.4cm,width=11cm]{imagenes/estudiantes}
	\end{center}
		
{\color{titulos} {\large \textbf{Descripción para la semana:}}}\\
		
Durante la semana, nos adentraremos en la exploración de métodos fundamentales para aproximar soluciones a ecuaciones diferenciales ordinarias de primer orden. Específicamente, dedicaremos tiempo a comprender en detalle el método de Euler y los métodos de Taylor. Estas técnicas proporcionan herramientas esenciales en el campo de la matemática aplicada y la física, permitiendo aproximaciones efectivas que nos acercan a comprender y resolver una variedad de fenómenos dinámicos en diferentes contextos. 
		
{\color{titulos} {\large \textbf{ Modo de trabajo:}}}
		
{ \fboxsep 12pt
	\fcolorbox {cuadro1}{cuadro1}{
		\begin{minipage}[l][10.5cm]{3.15cm}
			\begin{center}
				\includegraphics[scale=0.2]{imagenes/uno}
			\end{center}
			{\color{ContenidosObjetivos} {\large Material a utilizar}} \\ \\
			Libro de texto:\\
			$\bullet${\footnotesize 
			\ \ \URLLibroMN } \\			
			$\bullet${\footnotesize 
						\ \ \FolletoPractica}\\
			$\bullet${\footnotesize 
									\ \ \PaqueteSymbolic}\\
						\vspace{0.5cm}
						
		\end{minipage}
		}
}
{ \fboxsep 12pt
		\fcolorbox {cuadro2}{cuadro2}{
			\begin{minipage}[l][10.5cm]{3.15cm}
				\begin{center}
					\includegraphics[scale=0.2]{imagenes/dos}
				\end{center}
				{\color{ContenidosObjetivos} {\large Vídeos}} \\ \\
				{{ \tiny $\blacksquare$} \footnotesize   \URLVideoUno}\\
				{{ \tiny $\blacksquare$} \footnotesize   \URLVideoDos}\\
				{{ \tiny $\blacksquare$} \footnotesize   \URLVideoTres}\\ 
				{{ \tiny $\blacksquare$} \footnotesize   \URLVideoCuatro}\\ 
				{{ \tiny $\blacksquare$} \footnotesize Programación del Método de Euler en Octave}\\
				{{ $ ^{}$ \ \ \tiny $\blacksquare$} \footnotesize   \URLVideoCinco} \\
				{{ $ ^{}$ \ \ \tiny $\blacksquare$} \footnotesize   \URLVideoSeis}\\
				{{ \tiny $\blacksquare$} \footnotesize   \URLVideoSiete}\\ 
				{{ \tiny $\blacksquare$} \footnotesize   \URLVideoOcho}\\ 
%				{{ \tiny $\blacksquare$} \footnotesize   \URLVideoNueve}\\ 
				{{ \tiny $\blacksquare$} \footnotesize   \URLVideoDiez}\\ 
			\end{minipage}
	}
}
{ \fboxsep 12pt
	\fcolorbox {cuadro3}{cuadro3}{
		\begin{minipage}[l][10.5cm]{3.9cm}
			\begin{center}
				\includegraphics[scale=0.2]{imagenes/tres}
			\end{center}
			{\color{ContenidosObjetivos} {\large Evaluación}} \\ \\
			 \footnotesize\raggedright
\textbf{Clase presencial y quiz}\\
Mañana, \textbf{martes 27 de octubre}, la clase será presencial y se aplicará un quiz sobre \textbf{interpolación y regresión lineal}.

\medskip
\textbf{Tercer quiz virtual}\\
Esta semana se habilitará el tercer quiz virtual sobre \textbf{interpolación y regresión lineal}.

\medskip
\textbf{Defensas del laboratorio 1}\\
Se realizarán esta semana. Mañana, durante la clase, coordinaremos con los grupos las fechas y horarios.

\medskip
\textbf{Laboratorio 2}\\
Al finalizar la semana se asignará el segundo laboratorio, junto con las indicaciones para su desarrollo.
\end{minipage}
	} 
} 
\end{minipage}\hspace{0.2cm}\begin{minipage}[c]{5cm}
\begin{center}
{ \fboxsep 12pt
		\fcolorbox {CcontenidosObjetivos}{CcontenidosObjetivos}{
\begin{minipage}[l][21cm]{5cm}
\begin{center}
	\ImagenDM
\end{center}
					
{\color{ContenidosObjetivos}\textbf{ \small Contenidos:}}\\
					
{ \tiny $\blacksquare$} Ecuaciones diferenciales ordinarias\\ \\
{ \tiny $\blacksquare$} Método de Euler\\ \\
{ \tiny $\blacksquare$} Método de Taylor\\




{\color{white}\rule[0.1cm]{5cm}{0.01cm}}
		{\color{ContenidosObjetivos} \textbf{ \small Objetivos:}}\\
					
{ \tiny $\blacksquare$} Estudiar el método de Euler para la solución numérica de Ecuaciones diferenciales ordinarias de orden 1. \\ \\
{ \tiny $\blacksquare$} Programar el método de Euler. \\ \\
{ \tiny $\blacksquare$} Estudiar los métodos de Taylor para la solución numérica de Ecuaciones diferenciales ordinarias de orden 1. \\ \\
{ \tiny $\blacksquare$} Programar el método de Taylor de orden 2.



{\color{white}\rule[0.1cm]{5cm}{0.01cm}}\\

%\vspace{2cm}
%\hspace{-7pt}\begin{minipage}[4cm]{5cm}
%\begin{center}
%\includegraphics[scale=0.14]{imagenes/QRFacebook}
%
%Comunidad de la cátedra de Matemática Discreta en el tecDigital (\URLFacebook).
%\end{center}
%\end{minipage}

\end{minipage}
}}
\end{center}
\end{minipage}	 
%----------------------------------------------------------------------------


\newpage

%Página2----------------------------------------------------------------------
\begin{multicols}{2}
Se recomienda organizar su tiempo de manera que pueda trabajar por \textbf{etapas} con el fin de cubrir todos los contenidos y actividades, recuerde que el objetivo primordial es aprender. 
	 

Realice lo siguiente:
	
{\color{actividades} \textbf{  $\bullet$ Actividad 1: Ecuaciones Diferenciales. Introducción}}

Realice lo siguiente:
\begin{itemize}
\item Visualice los videos
\begin{itemize}
\item \URLVideoUno.
\item \URLVideoDos.
\end{itemize}
\item Del libro de texto, favor leer las páginas 376 a la 382. Hacer esta lectura de forma corrida, poniendo espacial atención en:
\begin{itemize}
\item ¿Qué es una Ecuación Diferencial?
\item ¿Qué es la solución general, una solución particular y una solución singular de una Ecuación diferencial?
\item ¿Qué es un problema de valor inicial?
\end{itemize}
\end{itemize}

{\color{actividades} \textbf{  $\bullet$ Actividad 2: Método de Euler}}

Realice lo siguiente:
\begin{itemize}
\item Visualice los videos:
\begin{itemize}
\item \URLVideoTres
\item \URLVideoCuatro
\end{itemize}
\item Del libro de texto lea la introducción: página 382-391. Estudie los ejemplos:
\begin{itemize}
\item Página 383, Ejemplos 138. 
\item Página 385: Ejemplo 385 y 140. 
\item Página 388: Ejemplo 141. 
\item Página 389: Ejemplo 142.

\end{itemize}
\item Del Folleto de prácticas estudiar:
\begin{itemize}
\item Página 86: Ejemplo 7.1. 
\item Página 87: Ejemplo 7.2
\end{itemize}
\end{itemize}

{\color{actividades} \textbf{  $\bullet$ Actividad 3: Ejercicios propuestos}}
\begin{itemize}
\item Ejercicios asignados:
\begin{itemize}
\item Del libro de texto. Realice todos los ejercicios del bloque (ejercicios 6.1) de la página 390.
\end{itemize}
\item Proceda a programar el métodos de Euler en octave, para lo cual se puede guiar de los siguientes videos:
\begin{itemize}
\item Programación del método de Euler \URLVideoCinco
\item Programación del método de Euler \URLVideoSeis
\end{itemize}
\end{itemize}

{\color{actividades} \textbf{  $\bullet$ Actividad 4: Métodos de Taylor}}

Realice lo siguiente:
\begin{itemize}
\item Visualice los videos:
\begin{itemize}
\item \URLVideoSiete
\item \URLVideoOcho
\end{itemize}
\item Del libro de texto lea la introducción: página 392-396. Estudie los ejemplos:
\begin{itemize}
\item Página 393: Ejemplos 143. 
\end{itemize}
\item Del Folleto de prácticas estudiar:
\begin{itemize}
\item Página 88: Ejemplo 7.3. 
\item Página 89: Ejemplo 7.4. Analice bien todos los aspectos mencionados, leer el documento las veces que sea necesario para entender totalmente lo que se plantea en ellos
\end{itemize}
\end{itemize}

{\color{actividades} \textbf{  $\bullet$ Actividad 5: Ejercicios propuestos}}
\begin{itemize}
\item Ejercicios asignados:
\begin{itemize}
\item Del libro de texto. Realice todos los ejercicios del bloque (ejercicios 6.2) de la página 395.
\end{itemize}
%\item Proceda a programar el métodos de Euler en octave, para lo cual se puede guiar de los siguientes videos:
%\begin{itemize}
%\item Programación del método de Euler \URLVideoNueve.
%\end{itemize}
\end{itemize}



%{\color{actividades} \textbf{  $\bullet$ Actividad 5: Métodos de Euler Mejorado}}
%
%Realice lo siguiente:
%\begin{itemize}
%\item Visualice los videos:
%\begin{itemize}
%\item \URLVideoDiez
%\end{itemize}
%\item Del libro de texto lea la introducción: página 396-399. Estudie el ejemplo:
%\begin{itemize}
%\item Página 397: Ejemplos 144. 
%\end{itemize}
%\item Del Folleto de prácticas estudiar:
%\begin{itemize}
%\item Página 88: Ejemplo 7.3. 
%\item Página 89: Ejemplo 7.4. Analice bien todos los aspectos mencionados, leer el documento las veces que sea necesario para entender totalmente lo que se plantea en ellos
%\end{itemize}
%\end{itemize}



\end{multicols}


{ \fboxsep 12pt
\fcolorbox {CRecapitulacion}{CRecapitulacion}{
\begin{minipage}[c]{12cm}
{\color{ContenidosObjetivos}\large {\large Recapitulación}:}\\
		
Se les recuerda la importancia de completar las actividades asignadas de la sección virtual para un desarrollo adecuado de su aprendizaje. El estudiante deberá realizar todas las actividades antes de la o las clases sincrónicas según lo indique el docente. \\
		
Se le recuerda que no está solo en el proceso. Puede realizar sus dudas en el grupo de Teams.\\

%Recuerde que esta semana se aplicará un quiz presencial, el cual ser realizará durante la clase del jueves 9 de mayo. Ver la sección de evaluaciones para más detalles. \\
		

\end{minipage}
		\begin{minipage}[c]{7cm}
			\begin{center}
				\Motivacion
			\end{center}
			
		\end{minipage}
	}
}

%{\color{actividades} \textbf{   $\bullet$ Códigos QR de los vídeos:}}\\
%\begin{center}
%\begin{tabular}{|c|c|c|c|}
%\hline 
%\textbf{¿Para que sirven las matemáticas?} & \textbf{Motivación a la lógica} & \textbf{Vídeo de teoría 1} &\textbf{Vídeo de teoría 2} \\ \hline &&&\\  
%\includegraphics[scale=0.1]{imagenes/QRMoti1} & \includegraphics[scale=0.1]{imagenes/QRMoti2} & \includegraphics[scale=0.1]{imagenes/QRVideo1} & \includegraphics[scale=0.1]{imagenes/QRVideo2}\\ 
%\hline 
%\textbf{Vídeo 3}&\textbf{Vídeo 4}&\textbf{Vídeo 5}&\\ 
%\hline &&&\\  
% \includegraphics[scale=0.1]{imagenes/QRVideo3} & \includegraphics[scale=0.1]{imagenes/QRVideo4} & \includegraphics[scale=0.1]{imagenes/QRVideo5} & \\ \hline
%\end{tabular} 
%\end{center}
\end{document}
